\chapter{Low-Rank Density Filters, KR Maps, and Retained Objects}
\label{ch:bf-highdim-low-rank-density-kr}
\label{ch:bf-highdim-particle-transport-tensor}
\label{prop:bf-hd-corrected-proposal}

\section{Orientation And Scope}
\label{sec:bf-hd-ttkr-purpose}

This chapter studies a high-dimensional nonlinear filtering problem through the
lens of one specific computational object: a low-rank non-Gaussian retained
representation built from tensor-train density approximations, conditional
Knothe--Rosenblatt maps, and squared-density repair.  The exact target remains
the filtering law and normalizer from
Chapter~\ref{ch:bf-highdim-nonlinear-foundations}; the approximation here is the
retained low-rank density or retained-object representation that numerically
replaces that law.

The scientific tension is the same one that appears across many high-dimensional
filtering problems.  The exact Bayesian recursion propagates a full conditional
law, but once the model is nonlinear that law is usually not available in a
finite-dimensional closed form.  Zhao and Cui answer that tension with a
non-Gaussian density-approximation lane based on tensor trains,
squared-density repair, conditional KR maps, and preconditioning.  The purpose
of this chapter is to make that retained-object construction readable in source
order and to state clearly what object is being carried, normalized, updated,
and later handed to the shared fixed-branch derivative chapter.

\paragraph{What this chapter does not claim.}
This chapter does not claim that adaptive TT/KR algorithms are globally smooth,
that moderate TT rank is always sufficient in high dimension, or that a retained
low-rank object automatically yields an exact posterior target for HMC.  Its
purpose is narrower: define the retained object, explain when that object is
scientifically meaningful, and state the diagnostics that later chapters must use
before promotion.

\section{Running Example And Notation}
\label{sec:bf-hd-ttkr-notation}

The running nonlinear example should stay conceptually continuous with the rest
of the block, but the retained-object lane needs a slightly richer notation than
the deterministic Gaussian lane.  The key distinction is between:
\begin{itemize}[leftmargin=0.28in]
  \item the exact Bayesian target law and its normalizer,
  \item the adaptive low-rank approximation family that numerically replaces it,
  \item and the fixed-branch scalar/derivative objects that later chapters may
  define from a frozen member of that family.
\end{itemize}
The method uses state variables \(x_t\), observation variables \(y_t\), and a
static parameter \(\alpha\).  Coordinate group notation such as
\(x_{t,<j},x_{t,>j},x_{t,\le j},x_{t,\ge j}\) matters because conditional
transport and conditional density construction are coordinate-local operations.
For TT/KR methods, the essential bookkeeping objects are coordinate groups,
conditional densities, retained marginal/filter objects, and any Jacobian terms
needed when the retained object is moved into preconditioned coordinates.

\section{Exact Filtering Target And Recursive Bottleneck}
\label{sec:bf-hd-ttkr-exact-target}

The retained-object lane starts from the same target discipline as the rest of
the block.  At time \(t\), the filtering problem is to update
\[
  p(x_t,\alpha\mid y_{1:t}),
\]
and the path problem is to update
\[
  p(x_{0:t},\alpha\mid y_{1:t}).
\]
After observing \(y_t\), the exact adjacent-state target has the three-block
form
\[
  q_t(x_t,\alpha,x_{t-1})
  =
  p(x_{t-1},\alpha\mid y_{1:t-1})
  f_\alpha(x_t\mid x_{t-1})
  g_\alpha(y_t\mid x_t),
\]
and the next filter is obtained by integrating out \(x_{t-1}\) and normalizing.
The bottleneck is therefore always the same: how to represent this object and
how to integrate the represented object without returning to an infeasible full
grid.

Implementation object.  A minimal implementation must expose evaluators for the
previous retained filter, the transition density, and the observation density,
then build an evaluator for the three-block target.  The first failure check is
that this target must be finite and nonnegative at all fitting points; otherwise
neither a square root nor an importance weight is meaningful.

\section{Tensor-Train Approximation Toolkit}
\label{sec:bf-hd-tt-toolkit}

A tensor train is a compressed representation of a multidimensional array or
function.  For a grid tensor \(A(i_1,\ldots,i_D)\), the TT format writes
\begin{equation}
  A(i_1,\ldots,i_D)=G_1(i_1)G_2(i_2)\cdots G_D(i_D),
  \label{eq:bf-hd-tt-format}
\end{equation}
where the \(G_k(i_k)\) are small matrices with connecting ranks.  If the mode
size is \(M\) and ranks are bounded by \(r\), storage is \(O(DMr^2)\) instead
of \(O(M^D)\).  The key condition is rank stability: if ranks grow with the
effective dimension, the compression disappears.

The chapter uses this toolkit for filtering objects, not for abstract storage
arithmetic alone.  The object being approximated may be a density, an
unnormalized adjacent-state target, a retained marginal, or a conditional map.
The approximation only helps if the represented object remains interpretable as a
probability or retained filtering quantity after truncation, marginalization, and
normalization.

\subsection{Rank-Count Plausibility}
\label{sec:bf-hd-rank-plausibility}

The method is plausible when the posterior geometry has limited interaction width
after ordering and preconditioning.  If every coordinate uses \(p\) basis
functions and the function has \(D\) coordinates, a full tensor stores
\(p^D\) coefficients, while a TT with moderate ranks stores approximately
\(DpR^2\) coefficients.  This is the theorem-to-implementation bridge: low-rank
structure is not decoration, but the empirical condition under which the method
escapes the full-grid curse of dimensionality.

Rank diagnostics therefore belong to the method contract.  If fitted ranks grow
persistently toward the cap or if conditioning deteriorates at the same time,
then the chosen ordering, basis, or preconditioner is not expressing the target
compactly.

\section{Coordinate Systems, Conditional Maps, And Retained Objects}
\label{sec:bf-hd-ttkr-coordinates}

The retained-object lane uses several coordinate systems because different
operations are easiest in different spaces.  Let
\begin{equation}
  r=(x_t,\alpha,x_{t-1})
  \label{eq:bf-hd-coord-physical}
\end{equation}
denote physical filtering coordinates.  A domain map sends a reference cube
coordinate \(z\in[-1,1]^D\) to physical coordinates:
\begin{equation}
  r=\Psi_t(z),
  \qquad
  J_{\Psi,t}(z)=\left|\det\nabla\Psi_t(z)\right|.
  \label{eq:bf-hd-coord-domain}
\end{equation}
The target seen by the TT in reference coordinates is
\begin{equation}
  \widetilde q_t(z)=q_t(\Psi_t(z))J_{\Psi,t}(z).
  \label{eq:bf-hd-coord-reference-density}
\end{equation}
Thus \(z\) is the coordinate in which basis functions, fitting points, and mass
matrices are declared.

Preconditioning introduces another coordinate \(u\).  Let \(T_t\) be an
invertible map from physical coordinates to a preconditioned coordinate:
\begin{equation}
  u=T_t(r),
  \qquad
  r=T_t^{-1}(u).
  \label{eq:bf-hd-coord-preconditioner}
\end{equation}
If \(\eta(u)\) is the reference density used in \(u\)-space, then
\begin{equation}
  q_t(r)\dd r
  =
  q_t(T_t^{-1}(u))\left|\det\nabla T_t^{-1}(u)\right|\dd u.
  \label{eq:bf-hd-coord-density-u}
\end{equation}
The retained coordinate \(z_t\) is the part of the reference coordinate retained
after marginalizing the previous state:
\begin{equation}
  z=(z_t,z_{t-1}),
  \qquad
  a_t(z_t)=\int \widehat q_t(z_t,z_{t-1})\dd z_{t-1},
  \label{eq:bf-hd-coord-retained}
\end{equation}
and
\begin{equation}
  \widehat p_t(z_t)=\frac{a_t(z_t)}{\widehat Z_t}.
  \label{eq:bf-hd-coord-retained-filter}
\end{equation}
Most implementation mistakes in this lane are coordinate mistakes: using a
physical density in reference coordinates without the Jacobian, applying a KR map
in \(u\)-space to a point still in \(r\)-space, or carrying a retained filter in
\(z_t\) while evaluating it as if it were already a physical density.

\subsection{One Observation Step Through All Coordinate Systems}
\label{sec:bf-hd-ttkr-coordinate-walkthrough}

The symbols \(r,z,u,z_t\) become easier to track when one observation update is
written as a single chain.  Begin in physical coordinates
\begin{equation}
  r=(x_t,\alpha,x_{t-1}),
  \qquad
  q_t(r)=
  \widehat p_{t-1}(x_{t-1},\alpha)
  f_\alpha(x_t\mid x_{t-1})
  g_\alpha(y_t\mid x_t).
  \label{eq:bf-hd-ttkr-walk-physical}
\end{equation}
A domain map
\begin{equation}
  r=\Psi_t(z),
  \qquad
  z\in[-1,1]^D,
  \qquad
  J_{\Psi,t}(z)=|\det\nabla\Psi_t(z)|
  \label{eq:bf-hd-ttkr-walk-domain}
\end{equation}
turns the physical density into a reference-coordinate density
\begin{equation}
  \widetilde q_t(z)=q_t(\Psi_t(z))J_{\Psi,t}(z).
  \label{eq:bf-hd-ttkr-walk-reference}
\end{equation}
The squared-TT route fits a source-scale square root in reference coordinates and
then returns to the unshifted evidence scale through
\begin{equation}
  q_{t,{\rm src}}(z)=\phi_t(z)^2+\tau_t\lambda_t(z),
  \qquad
  \widehat q_t(z)=e^{-c_t}q_{t,{\rm src}}(z).
  \label{eq:bf-hd-ttkr-walk-source-evidence}
\end{equation}
The carried filter is not the full joint object.  If \(z=(z_t,z_{t-1})\), then
\begin{equation}
  a_t(z_t)=\int \widehat q_t(z_t,z_{t-1})\dd z_{t-1},
  \qquad
  \widehat p_t^{\rm ref}(z_t)=\frac{a_t(z_t)}{\widehat Z_t}.
  \label{eq:bf-hd-ttkr-walk-retained}
\end{equation}
If a later formula needs the filter as a physical density, write
\begin{equation}
  (x_t,\alpha)=\Psi_{t,{\rm ret}}(z_t),
  \qquad
  \widehat p_t^{\rm phys}(x_t,\alpha)=
  \widehat p_t^{\rm ref}(\Psi_{t,{\rm ret}}^{-1}(x_t,\alpha))
  \left|\det\nabla\Psi_{t,{\rm ret}}^{-1}(x_t,\alpha)\right|.
  \label{eq:bf-hd-ttkr-walk-physical-filter}
\end{equation}
The bookkeeping rule is: \(\widetilde q_t(z)\), \(a_t(z_t)\), \(\widehat Z_t\),
and \(\widehat p_t^{\rm ref}(z_t)\) are reference-coordinate objects integrated
with \(\dd z\); \(\widehat p_t^{\rm phys}\) is the corresponding physical
object after the declared Jacobian correction.

Preconditioning inserts one additional coordinate \(u\).  A bridge density
\(\rho_t(r)\) and map \(T_t:r\mapsto u\) are chosen so that the dominant shape
of \(q_t\) is easier in \(u\)-space.  The residual target is
\begin{equation}
  q_t^\sharp(u)=
  \frac{q_t(T_t^{-1}(u))}{\widehat\rho_t(T_t^{-1}(u))}\eta(u),
  \label{eq:bf-hd-ttkr-walk-residual}
\end{equation}
and the pullback relation after fitting is
\begin{equation}
  \widehat q_t(r)=
  \frac{\widehat\nu_t^\sharp(T_t(r))}{\eta(T_t(r))}\widehat\rho_t(r).
  \label{eq:bf-hd-ttkr-walk-pullback}
\end{equation}
Thus the full coordinate story of one observation step is
\begin{equation}
  q_t(r)
  \to
  \widetilde q_t(z)
  \to
  \widehat q_t(z)
  \to
  \widehat p_t^{\rm ref}(z_t)
  \to
  \widetilde q_{t+1}(z_{t+1},z_t),
  \label{eq:bf-hd-ttkr-walk-full-chain}
\end{equation}
with the preconditioned route adding
\begin{equation}
  q_t(r)\to q_t^\sharp(u)\to\widehat\nu_t^\sharp(u)\to\widehat q_t(r).
  \label{eq:bf-hd-ttkr-walk-preconditioned-chain}
\end{equation}
The filtering task does not change; only the coordinate system in which the hard
part of the density is approximated changes.

\section{TT Sequential Learning And Conditional KR Transports}
\label{sec:bf-hd-zhao-cui}

\citet{zhao2024ttsequential} is the central discrete-time TT source for this
lane.  Its checked role here is direct: the chapter uses the paper’s sequential
state-and-parameter learning setup, its posterior-density approximation view,
and its conditional Knothe--Rosenblatt map construction from a TT approximation.

A nonnegative TT approximation supplies conditional density factors or
conditional cumulative distribution functions.  A triangular conditional map can
then be constructed from one-dimensional conditional CDFs, which is the KR
transport idea in operational form.  For BayesFilter this is attractive because
it connects low-rank density approximation to transport preconditioning.  It
remains conditional because normalization, monotonicity, rank error, and
Jacobian stability are part of the method, not afterthoughts.

\section{What This Chapter Exports Forward}
\label{sec:bf-hd-ttkr-front-half-exports}

At the end of this front-half chapter, the retained-object lane has fixed the
objects that later operational chapters are allowed to manipulate:
\begin{enumerate}[leftmargin=0.28in]
  \item the adjacent-state target and its coordinate systems,
  \item the TT approximation toolkit and rank-count plausibility logic,
  \item the retained numerator \(a_t\) and normalized carried filter
  \(\widehat p_t\),
  \item the KR / preconditioning map perspective,
  \item the core diagnostics that decide whether the retained object remains
  scientifically meaningful.
\end{enumerate}
The next chapter,
Chapter~\ref{ch:bf-highdim-squared-tt-fixed-branch-likelihoods}, then takes up
the operational middle: squared-TT recursion, fixed-branch likelihood
construction, stored fields, and the exact scalar object exported to the
same-scalar derivative chapter.

\section{Tensor-Network Kalman And Square-Root Caution}
\label{sec:bf-hd-tn-kalman}

Tensor-network covariance compression is included here only as a cautionary
adjacent lane.  If a covariance-like object is compressed, the method must expose
either a PSD reconstruction check or a factor/square-root representation whose
product is positive semidefinite by construction.  This material remains
secondary to the retained TT/KR density lane; it is kept because it clarifies one
important boundary between density compression and covariance compression.

\section{Complexity, Diagnostics, And Exports}
\label{sec:bf-hd-ttkr-complexity}

The retained-object lane exports five primary diagnostics:
\begin{enumerate}[label=(\arabic*), leftmargin=0.28in]
  \item rank growth,
  \item normalization/mass error,
  \item positivity / nonnegativity construction,
  \item support / Jacobian / map validity,
  \item retained-object consistency for the next time step.
\end{enumerate}
These diagnostics are not side comments; they are the mechanism by which the
chapter decides whether the retained object remains scientifically meaningful.

\paragraph{Exported forward.}
The next operational chapter may import only the retained-object scalar
ingredients and branch-defining structural choices that this chapter makes
explicit.  The later validation/promotions chapter may import only the defects,
diagnostics, promotion gates, and source-risk boundaries stated here.  Neither
later chapter should silently strengthen TT/KR transport existence,
branch-local differentiability, or retained-object feasibility beyond what the
checked sources support.

\section{Methodological Boundary And Sources}
\label{sec:bf-hd-ptt-boundary}

The substantive claims in this chapter are tied to TT, KR, transport, and
adjacent tensor-network sources where their checked constructions are used.
Savostyanov maxvol theory and Knothe's original rearrangement source were not
inspected directly here; the chapter therefore uses Oseledets--Tyrtyshnikov for
TT-cross and Rosenblatt plus modern transport-filtering papers for triangular
maps.  No source here is used to claim production readiness, NAWM readiness,
posterior accuracy, tensor-method validation, transport-method validation, HMC
convergence, GPU/XLA readiness, or default readiness.
